\documentclass[11pt]{exam} 
\usepackage[lmargin=1in,rmargin=1.75in,bmargin=1in,tmargin=1in]{geometry}  


\input{preamble}



\begin{document}
	\week{10: Randomized Algorithms}{April 7, 2026}
	%\lecture{19: Fingerprinting}{April 7}
	%\lecture{20: Hashing}{April 9}
	%\lecture{21: PageRank}{April 14}
	
	
	\paragraph{Course Logistics}
	\begin{itemize}
		\item Reading from Chapters 5 and 11 this week, especially sections on randomized algorithms and hashing
		\item Homework 6 due Friday
	\end{itemize}
	
	\section{The Matrix Verification and Fingerprinting Problem}
	Consider the problem of checking if two large structures are identical. A deterministic approach might require reading all elements, taking a long time. 
    Instead, we can use \emph{randomization} to map the large structure to a smaller \emph{fingerprint}.\\
	
	Given two $n \times n$ matrices $A, B$, and a proposed product $C$, checking if $A \times B = C$ deterministically takes at least $O(n^\omega)$ time.
	
	\vs{2.5cm}
	
	
	Using a random $n \times 1$ vector $r \in \{0,1\}^n$, the randomized fingerprint check is given by: 
\vs{3cm}
	
	Let $\mathcal{F}$ be the set of valid runs. The \emph{fingerprint comparison} is:
	\begin{equation*}
		\text{match}(A, B, C, r) = \begin{cases}
		 \\  \\%	\text{True} & \text{if } A \times (B \times r) = C \times r \\ \\
			\text{False} & % \text{otherwise}
		\end{cases}
	\end{equation*}
	A \hide{false positive} occurs when $A \times B \neq C$ but the algorithm still outputs True. \\
	
	\paragraph{Definition} A \emph{Monte Carlo} algorithm is a randomized algorithm whose runtime is deterministic, but whose output has a bounded probability of being incorrect.
	
	\newpage
	
	
	\subsection{Fingerprinting Properties}
	
	\paragraph{Probability of Error}
	For any matrices $A$, $B$, and $C$ where $A \times B \neq C$, if $r \in \{0,1\}^n$ is chosen uniformly at random,
	\begin{equation*}
		\text{Pr}[A(Br) = Cr] \leq \frac{1}{2} \\
	\end{equation*}
	
	\paragraph{Optimal Substructure for Evaluation}
	\begin{lemma} If $r$ is an $n \times 1$ vector, the matrix-vector multiplications $B \times r$ and then $A \times (Br)$ \hide{take $O(n^2)$ time.}
	\end{lemma}
	
	\vs{5cm}
	
	\paragraph{Las Vegas Algorithms}
	\begin{lemma}: In contrast to Monte Carlo algorithms, a Las Vegas algorithm %always returns the correct answer, but its runtime is a random variable.
	\end{lemma}
	\vs{2cm}
	
	
	\newpage
	
	\subsection{Error Amplification}
	\begin{Qu}
		Assume we run Freivalds' algorithm $k$ independent times to check if $A \times B = C$. If the algorithm returns True every time, what is the probability that $A \times B \neq C$ (i.e., we are seeing a false positive)?
		\begin{itemize}
			\aitem It will always be exactly $1/2$
			\bitem It will be at most $1/k$
			\citem It will be at most $(1/2)^k$
			\ditem The probability will drop to exactly 0 after $k=2$ runs
		\end{itemize}
	\end{Qu}
	\vs{5cm}
	
	\newpage
	\subsection{Freivalds' Algorithm Basics}
	Let $k$ denote the number of independent trials. In our algorithm, we maintain the following attributes for each trial $i$: \\
	
	
	\begin{tabular}{| l | p{4.5cm} | p{4.5cm} | p{4.5cm} |}
		\hline
		\textbf{Attribute} & \textbf{Explanation} & \textbf{Initialization} & \textbf{Invariant} \\
		\hline
		$trial.r$ & \phantom{the random vector chosen for this iteration} \newline &  \phantom{$r \sim \{0,1\}^n$;} \newline \phantom{randomize} & \phantom{We will always maintain:  $r_j \in \{0,1\}$} \\
		\hline
		$trial.result$ \newline &  \phantom{whether the fingerprint matched} \newline \phantom{whether the fingerprint matched} &  & \\
		\hline
	\end{tabular}
	
	\vs{2cm}
	
	Given values for these attributes, we can construct an overall \emph{decision} structure:
	\begin{itemize}
		\item If any trial returns False, the matrices are definitely not equal. \\\
		\item If all trials return True, we assert they are equal with high confidence.
	\end{itemize}
	
	
	If $trial.result = \text{True}$ for all $k$ trials, then the confidence graph $G_c$ is called a \hide{high-probability match}.
	
	\vs{3cm}
	
	
	
	\newpage
	
	\subsection{Generic Fingerprinting Algorithm}
	The key idea behind randomized fingerprinting is to successively ``test'' randomly generated fingerprints.
	\begin{algorithm}
		\textsc{TestFingerprint}($A,B,C$)
		\begin{algorithmic}
			\State Generate a random vector $r \in \{0,1\}^n$
			\If{$A \times (B \times r) \neq C \times r$}
			\State 
			\State
			\State
			\State
			\EndIf
		\end{algorithmic}
	\end{algorithm}
	
	\vs{5cm}
	\begin{algorithm}
		\textsc{GenericFreivalds}($A,B,C,k$)
		\begin{algorithmic}
			\State \textsc{Initialize-Trials}$(k)$
			\While{We have not exceeded $k$ trials}
			\State 
			\State
			\State
			\State
			\EndWhile
		\end{algorithmic}
	\end{algorithm}
	
	\newpage
	\includegraphics[width = .45\linewidth]{small-graph-fingerprint.png}
	
	\includegraphics[width = .45\linewidth]{small-graph-fingerprint.png}
	
	\includegraphics[width = .45\linewidth]{small-graph-fingerprint.png}
	
	\includegraphics[width = .45\linewidth]{small-graph-fingerprint.png}
	
	\newpage 
	
	\section{The Universal Hashing Algorithm}
	
	Let's do a quick recap of what we have learned about dictionaries and hashing.

	\paragraph{Adversarial Inputs} Let $h$ be a hash function mapping keys to a table of size $m$.
	
	\begin{itemize}
		\item Case 1: $h$ is fixed and deterministic: an adversary can choose $n$ keys that all hash to the same slot, leading to $O(n)$ search time. 
		\item Case 2: $h$ is chosen randomly from a universal family: then the expected number of collisions is bounded, and we have $O(1)$ expected search time. 
	\end{itemize}
	
	If Case 1 is true, we would prefer to be aware of this, in which case we state that the worst-case time is $O(n)$ and we're done.
	
	Otherwise, we use randomized hash families when solving the dictionary problem.
	
	\paragraph{Hashing and Generic Algorithm} 
	\begin{algorithm}
		\textsc{HashInsert}($T, x, h$)
		\begin{algorithmic}
			\If{$T[h(x)]$ is empty}
			\State $T[h(x)] = x$
			\Else
			\State Append $x$ to the chain at $T[h(x)]$
			\EndIf
		\end{algorithmic}
	\end{algorithm}
	\begin{algorithm}
		\textsc{GenericHashing}($S, m$)
		\begin{algorithmic}
			\State \textsc{Initialize-Hash-Table}$(m)$
			\State Choose $h \in \mathcal{H}$ uniformly at random
			\While{There are elements $x \in S$ to insert}
			\State \textsc{HashInsert}($T, x, h$)
			\EndWhile
		\end{algorithmic}
	\end{algorithm}
	
	Now we are ready to find a more concrete implementation of this approach, by answering the question:
	

	\newpage
	\subsection{Universal Hashing Algorithm Pseudocode}
	
	\begin{algorithm}
		\textsc{Universal-Hash-Table}($S, m, p$)
		\begin{algorithmic}
			\For{$i = 0$ to $m - 1$}
			\State $T[i] = \text{NIL}$
			\EndFor
			\State $a = \textsc{Random}(1, p-1)$
			\State $b = \textsc{Random}(0, p-1)$
			\State 
			\For{$ x \in S$}
			\State $h(x) = ((a \cdot x + b) \bmod p) \bmod m$
			\State $\textsc{HashInsert}(T, x, h)$
			\EndFor
			\State 
			\For{$ i = 0 \text{ to } m-1$}
			\If{$\text{Length}(T[i]) > 1 + \frac{|S|}{m}$}
			\State Return ``the chain is longer than expected"
			\EndIf
			\EndFor
			\State
			\State $V_H = \{ x \in S \}$
			\State $E_H = \{(x, y) \colon h(x) = h(y)\}$
			\State Return $T$ and $\{ \text{Length}(T[i]) \colon i \in \{0 \dots m-1\}\}$
		\end{algorithmic}
	\end{algorithm}
	
	\begin{Qu}
		What is the expected runtime complexity of an unsuccessful search using universal hashing with chaining, where $\alpha = n/m$ is the load factor?
		\begin{itemize}
			\aitem $O(1 + \alpha)$
			\bitem $O(\log n)$
			\citem $O(n / m + \log m)$
			\ditem $O(n)$
		\end{itemize}
	\end{Qu}
	
	\newpage
	
	\begin{algorithm}[t]
		\textsc{Universal-Hashing}($S,m,p$)
		\begin{algorithmic}
			\State \textsc{InitializeHashTable}($m$)
			\State $a = \textsc{Random}(1, p-1)$, $b = \textsc{Random}(0, p-1)$
			\For{$ x \in S$}
			\State $h(x) = ((a \cdot x + b) \bmod p) \bmod m$
			\State $\textsc{HashInsert}(T, x, h)$
			\EndFor
			\State Return $T$
		\end{algorithmic}
	\end{algorithm}
	\includegraphics[width = \linewidth]{twohash.png}
	
	\includegraphics[width = \linewidth]{twohash.png}
	
	\includegraphics[width = \linewidth]{twohash.png}
	
	\includegraphics[width = \linewidth]{twohash.png}
	
	\newpage
	
	\subsection{Correctness}
	\begin{theorem}
		Assuming that the hash family $\mathcal{H}_{p,m}$ is universal, when \textsc{Universal-Hashing} is used, the expected length of the chain $E[\text{Length}(T[h(x)])]$ will equal $1 + \alpha$ for every $x \in S$. 
	\end{theorem}
	
	
	\newpage
	
	
	\subsection{Catching Bad Hash Functions}
	\begin{algorithm}
		\textsc{Rehash-If-Bad} (Collision check)
		\begin{algorithmic}
			\State Do everything else...
			\For{$ i = 0 \text{ to } m-1$}
			\If{$\text{Length}(T[i]) > 2 \cdot (1 + \alpha) \cdot \log n$}
			\State Return an error, the hash function chosen was poor
			\EndIf
			\EndFor
		\end{algorithmic}
	\end{algorithm}
	
	\begin{theorem}
		If we pick $h$ randomly from a universal family, the probability that the code above produces an error and we need to rehash is very small (by Markov's/Chernoff bounds).
	\end{theorem}
	
	\newpage
	
		\section{PageRank Algorithm}
	\begin{algorithm}
		\textsc{PageRank}$(G, d, \text{tol})$
		\begin{algorithmic}
			\State \textsc{Initialize-PageRank}$(G, d)$
			\State $S = \emptyset$
			\State $Q = V$
			\While{$\text{diff} > \text{tol}$}
			\State $u = \textsc{ExtractMaxChange}(Q)$
			\State $S = S \cup \{u\}$
			\For{$v \in \text{Adj}[u]$}
			\State \textsc{UpdateRank}($u,v,d$)
			\EndFor
			\EndWhile
		\end{algorithmic}
	\end{algorithm}
	
	This algorithm maintains:
	\begin{itemize}
		\item $PR[u]$ = current estimate of the stationary probability for node $u$
		\item $d$ = damping factor (teleportation probability) to handle spider traps and dead ends
	\end{itemize}
	\vs{2cm}
	\begin{Qu}
		Why do we introduce the damping factor $d$ in the PageRank algorithm?
		\begin{itemize}
			\aitem To speed up the adjacency list traversal
			\bitem To guarantee convergence by making the transition matrix strongly connected and aperiodic
			\citem To eliminate all nodes with in-degree 0
			\ditem To reduce the memory usage of the algorithm
		\end{itemize}
	\end{Qu}
	
	
	\newpage
	
	\subsection{Example}
	\includegraphics[width = .45\linewidth]{small-graph-pagerank.png}
	
	\includegraphics[width = .45\linewidth]{small-graph-pagerank.png}
	
	\includegraphics[width = .45\linewidth]{small-graph-pagerank.png}
	
	\includegraphics[width = .45\linewidth]{small-graph-pagerank.png}
	\newpage
	\subsection{Correctness}
	\vspace{-20pt}
	\begin{algorithm}
		\textsc{PageRank-PowerIteration}$(G, d, \text{iters})$
		\begin{algorithmic}
			\State \textsc{Initialize-PageRank}$(G)$
			\For{$i = 1$ to $\text{iters}$}
			\State $\text{next\_PR} = \text{array of zeros}$
			\For{$u \in V$}
			\State $\text{next\_PR}[u] = \frac{1-d}{|V|}$
			\For{$v \in In(u)$}
			\State $\text{next\_PR}[u] += d \cdot \frac{PR[v]}{Out(v)}$
			\EndFor
			\EndFor
			\State $PR = \text{next\_PR}$
			\EndFor
		\end{algorithmic}
	\end{algorithm}
	We will take the following properties as given.
	
	\paragraph{Stochastic property}
	For each column in the modified transition matrix $M'$, the sum of the elements is exactly $1$.\\
	
	\paragraph{Convergence property}
	Let $\pi$ be the stationary distribution. By the Perron-Frobenius Theorem, iterative multiplication by a positive, column-stochastic matrix $M'$ will converge to the dominant eigenvector.\\
	
	\paragraph{Teleportation property}
	If there is a spider trap in the graph, the $1-d$ probability of teleporting ensures the random surfer does not get permanently stuck. \\ \\
	
	
	\begin{theorem}
		PageRank algorithm, run on a directed web graph $G = (V,E)$ with damping factor $d < 1$, terminates with $PR[u] = \pi(u)$ for all $u \in V$, where $\sum \pi(u) = 1$.
	\end{theorem}
	\textit{Proof.} We maintain the following invariant throughout. \\
	
	\textbf{Invariant}: At the start of the while loop, the vector $PR$ satisfies $\sum_{u \in V} PR[u] = 1$. \\
	
	If this invariant is true, the theorem is shown. Why?
	\newpage
	
	Assume this invariant is \emph{not} true and we will show a contradiction. \\
	
	Define $t$ to be the first iteration where the sum of probabilities $PR_t$ is not $1$.
	
	\begin{enumerate}
		\item We know $t > 0$, since $PR_0$ is initialized to $1/|V|$ for all $u$. \\
		\item We know $\sum_{u \in V} PR_{t-1}[u] = 1$. \\
		\item Let $M'$ be the modified transition matrix combining the web graph and teleportation. \\ 
		\vs{4cm} 
		\item We know that $PR_t = M' \times PR_{t-1}$\\
		\item We then know $\sum_u PR_t[u] = \sum_u \sum_v M'_{uv} PR_{t-1}[v]$ \\
		\item We know that this can be rewritten as $\sum_v \left( PR_{t-1}[v] \sum_u M'_{uv} \right)$ \\
		
		\item And we know $\sum_u M'_{uv} = 1$ because $M'$ is column-stochastic \\
		\item We know also that $\sum_v PR_{t-1}[v] = 1$ from our assumption \\
		\item The last steps show $\sum_u PR_t[u] = \sum_v PR_{t-1}[v] \cdot 1 = 1$. Why is this a contradiction?
	\end{enumerate}
	
	\newpage
	
	
	\subsection{Runtime Analysis}
	\begin{algorithmic}
		\State \textsc{Initialize-PageRank}$(G)$
		\For{$i = 1$ to $\text{iters}$}
		\State $\text{next\_PR} = \text{array of zeros}$
		\For{$u \in V$}
		\State $\text{next\_PR}[u] = \frac{1-d}{|V|}$
		\For{$v \in In(u)$}
		\State $\text{next\_PR}[u] += d \cdot \frac{PR[v]}{Out(v)}$
		\EndFor
		\EndFor
		\State $PR = \text{next\_PR}$
		\EndFor
	\end{algorithmic}
	\begin{Qu}
		Assuming the graph is stored using an adjacency list containing incoming edges for each node, what is the runtime of a single iteration of PageRank? (Warning: there may be more answer that you can argue is valid! Don't just select an answer, think about a specific reasoning for the runtime analysis you select).
		\begin{itemize}
			\aitem $O(V + E)$ 
			\bitem $O(V^2)$
			\citem $O(VE)$
			\ditem $O(E \log V)$
			\eitem $O(E + V \log V)$
		\end{itemize}
	\end{Qu}
	
	
	\newpage
	\subsection{Class Activity: another example}
	\includegraphics[width = .75\linewidth]{pagerankalg.png}
	
\end{document}